\documentclass[11pt]{article}
\usepackage[margin=1in]{geometry}
\usepackage{enumitem}
\usepackage{booktabs}
\usepackage{amsmath}
\usepackage[hidelinks]{hyperref}

\title{Temporal Diversification Weekly Playbook}
\author{Flyxion}
\date{2026}

\begin{document}
\maketitle

\section{Goal}

Sustain engagement by managing a portfolio of unfinished histories on different clocks,
not by demanding continuous novelty from a single process.

\section{Weekly Portfolio Construction}

Maintain at least four active scopes:
\begin{itemize}[leftmargin=*]
  \item \textbf{Fast} ($\tau\approx$ minutes--hours): daily execution and immediate feedback.
  \item \textbf{Medium} ($\tau\approx$ days): experiments and short deliverables.
  \item \textbf{Slow} ($\tau\approx$ weeks): proofs, architecture shifts, chapter drafts.
  \item \textbf{Very slow} ($\tau\approx$ months+): foundational programs and long manuscripts.
\end{itemize}

\section{Operational Rules}

\begin{enumerate}[leftmargin=*]
  \item \textbf{Never require all scopes to POP weekly.} Long horizons stay unresolved by design.
  \item \textbf{Require anti-starvation updates.} Every slow scope must receive at least one
  structural event per week.
  \item \textbf{Default action under local flattening: BIND.} If immediate closure is premature,
  add a constraint that sharpens next continuation.
  \item \textbf{Use REFUSE to cut dead branches quickly.} Remove known-nonproductive continuations.
  \item \textbf{Use COLLAPSE to reduce bookkeeping overhead.} Merge distinctions no longer
  decision-relevant.
\end{enumerate}

\section{Daily Scope-Selection Procedure}

At each work block:
\begin{enumerate}[leftmargin=*]
  \item Estimate candidate gains for active scopes: $\max_{a\in O_i}\mathrm{EMCG}_i(a\mid F_t)$.
  \item If a slow scope is near starvation threshold, prioritize it.
  \item Otherwise choose action by highest gain-per-cost.
  \item Log the event type (POP/REFUSE/BIND/COLLAPSE) and whether saturation or recovery was observed.
\end{enumerate}

\section{Weekly Review Metrics}

Track:
\begin{itemize}[leftmargin=*]
  \item time-to-local-saturation (per scope),
  \item total time in field saturation,
  \item recovery latency after switching scopes,
  \item starvation rate for slow scopes,
  \item stability of compression progress across horizons.
\end{itemize}

\section{Applied Templates}

\subsection*{Personal Workflow}
One fast execution scope, one medium experimentation scope, one slow integration scope,
one very-slow thesis scope. Weekly requirement: no slow scope receives zero events.

\subsection*{Education}
Design schedules that interleave immediate drills, multi-day projects, and long-form inquiry.
Measure whether boredom events correlate with local flattening in only one horizon.

\subsection*{Research and Programming}
Treat deep projects as unresolved scopes; require periodic BIND-like refinement events
(lemma sketch, boundary condition, interface constraint) even without POP.

\subsection*{AI Agent Orchestration}
Use multi-clock schedulers where long-horizon tasks get explicit anti-starvation allocations
instead of purely greedy short-horizon optimization.

\end{document}
